\documentclass[10pt,a4paper]{amsart}
\usepackage[margin=19mm]{geometry}
\usepackage{amsmath,amssymb,mathtools}
\usepackage[scr=rsfs,cal=euler]{mathalfa}
\usepackage{xfrac}
\usepackage[unicode,hidelinks,bookmarks=false]{hyperref}
\newcommand{\ie}{i.e.\ }
\newcommand{\cf}{cf.\ }
\newcommand{\resp}{respectively\ }
\newcommand{\Subsection}{\subsection}
\providecommand{\nobreakdash}{\nobreak\hbox{-}}
\setlength{\emergencystretch}{2em}
\multlinegap0pt
\usepackage{amssymb}
\usepackage{dsfont}
\usepackage{xcolor}
\usepackage[shortlabels]{enumitem}
\usepackage{mathtools}
\usepackage{tcolorbox}
\usepackage[normalem]{ulem}
\newtheorem{proposition}{Proposition}
\newcommand{\R}{{\mathbb R}}
\newcommand{\be}{\begin{equation}}
\newcommand{\ee}{\end{equation}}
\title{Potential Games on Unimodular Random Graphs}
\author{Eyal Neuman, Sturmius Tuschmann}
\date{Source: arXiv:2604.13836v1 (CC BY 4.0)}
\begin{document}
\maketitle

\noindent\emph{Source and licence.} Potential Games on Unimodular Random Graphs. Version arXiv:2604.13836v1 (CC BY 4.0), distributed by arXiv under the Creative Commons Attribution 4.0 International licence, http://creativecommons.org/licenses/by/4.0/, as recorded in the arXiv licence metadata for this exact version and shown on the abstract page https://arxiv.org/abs/2604.13836v1. Full licence notice: \texttt{licenses/arxiv-2604.13836-CC-BY-4.0.txt}.\par\smallskip

\noindent\textit{Editorial scope.} Counted unit: the article's proposition prop:W-example (source0.tex lines 683--690). The article's complete original proof of it is delivered with it (source0.tex lines 1750--1788, one \texttt{proof} environment of 39 lines, which the article places away from the statement). No mathematics is written, repaired or re-derived: the statement and the proof are the article's own lines, in the article's own notation, and no retained line is substituted or re-flowed.  Retained uncounted as the source-local context the proof needs: lines 218--222; lines 660--664; lines 666--670; lines 673--678.  Nothing was omitted from any retained span, so the package carries no omission record.  No retained line contains a citation, so the wrapper carries no bibliography and no bibliography entry.  No retained line contains a figure environment, an image-inclusion command or drawing code, so no image is delivered and the package declares no asset dependency.  Editorial formatting only, in addition: an AMS \texttt{amsart} preamble that restates the article's own declarations and macros used by the retained lines, namely the environments \texttt{proposition} on separate counters, together with the standard packages those lines need.  The retained environments print in this document as Proposition 1 in order of appearance, whereas the article prints the same items with the numbering of its own full document.  The complete native source of the article as distributed in the arXiv e-print for this exact version is bundled unchanged as \texttt{sources/198500/source0.tex}.
\begin{equation} \label{eq:J}
J_v^G(x^G)
:=
f([G,v],x^G_v)+\sum_{u\sim v} h([G,v,u],x^G_v,x^G_u).
\end{equation}

\begin{equation}\label{eq:W-f}
f:\mathcal G_\ast\times\mathbb R\to\mathbb R,
\quad
f([G,o],a):=-\eta([G,o])\,a,
\end{equation}

\begin{equation}\label{eq:W-h}
h:\mathcal G_{\ast\ast}\times\mathbb R\times\mathbb R\to\mathbb R,
\quad
h([G,o,v],a,b):=W(a-b),
\end{equation}

\begin{equation}\label{eq:W-Delta}
(\Delta_W x)([G,o])
:=
\frac1{\deg_G(o)}
\sum_{v\sim o}W'\!\big(x([G,o])-x([G,v])\big).
\end{equation}

\begin{proposition}\label{prop:W-example}
If \(x^\ast\in\mathcal H\) satisfies
\begin{equation}\label{eq:W-poisson}
(\Delta_W x^\ast)([G,o])=\frac{\eta([G,o])}{\deg_G(o)},
\quad \text{for } \mu\text{-a.e. } [G,o],
\end{equation}
then \(x^\ast\) is a quenched Nash equilibrium.
\end{proposition}

\begin{proof}[Proof of Proposition~\ref{prop:W-example}]
Let \(x^\ast\in\mathcal H\) satisfy \eqref{eq:W-poisson}. Recalling \eqref{eq:W-Delta}, fix a rooted graph \([G,o]\) in the full-measure set where \eqref{eq:W-poisson} holds, and define $\Gamma_{[G,o]}:\R\to\R$ by
\be\begin{aligned}\label{eq:W-Gamma}
\Gamma_{[G,o]}(a)
:&=
f([G,o],a)+\sum_{v\sim o}h([G,o,v],a,x^\ast([G,v]))
\\&=
-\eta([G,o])\,a+\sum_{v\sim o}W\!\big(a-x^\ast([G,v])\big),
\end{aligned}\ee
where we used \eqref{eq:W-f} and \eqref{eq:W-h}. Since \(W\) is convex, the map \(a\mapsto \Gamma_{[G,o]}(a)\) is convex. Moreover, by \eqref{eq:W-Gamma},
\[
\Gamma_{[G,o]}'(a)
=
-\eta([G,o])+\sum_{v\sim o}W'\!\big(a-x^\ast([G,v])\big).
\]
Using \eqref{eq:W-Delta} and \eqref{eq:W-poisson}, we obtain
\[
\Gamma_{[G,o]}'\big(x^\ast([G,o])\big)
=
-\eta([G,o])+\sum_{v\sim o}W'\!\big(x^\ast([G,o])-x^\ast([G,v])\big)
=
0.
\]
Therefore \(x^\ast([G,o])\) minimizes \(\Gamma_{[G,o]}\) over \(\mathbb R\). Moreover, by \eqref{eq:J} and \eqref{eq:W-Gamma}, for every \(a\in\mathbb R\),
\[
J_o^G(a,x_{-o}^{\ast,G})
=
f([G,o],a)+\sum_{v\sim o}h([G,o,v],a,x^\ast([G,v]))
=
\Gamma_{[G,o]}(a).
\]
Hence
\[
J_o^G(x^{\ast,G})=\Gamma_{[G,o]}(x^\ast([G,o]))
\le \Gamma_{[G,o]}(a)=J_o^G(a,x_{-o}^{\ast,G}),
\quad \text{for all } a\in\mathbb R.
\]
Since this holds for $\mu$-a.e.~\([G,o]\), \(x^\ast\) is a quenched Nash equilibrium.
\end{proof}

\end{document}
